\chapter{Trabalhos Relacionados} \label{trabs}

A literatura sobre \textit{sycophancy} em LLMs apresenta abordagens complementares: avaliações comportamentais, investigação da influência de preferências humanas, análise de diálogos com múltiplos turnos e distinção entre correções válidas e sugestões incorretas. Esta seção discute os quatro estudos selecionados como base para o TCC, destacando suas contribuições e os limites de sua aplicação ao recorte proposto.

\section{Descoberta do fenômeno}

\textcite{perez2023discovering} propuseram usar modelos de linguagem para gerar avaliações de diferentes comportamentos. No experimento sobre \textit{sycophancy}, modelos maiores apresentaram maior tendência a reproduzir a posição informada pelo usuário em perguntas de política, filosofia e pesquisa em processamento de linguagem natural. Entretanto, os autores relataram níveis semelhantes desse comportamento entre modelos com diferentes quantidades de etapas de RLHF (\textit{Reinforcement Learning from Human Feedback}), incluindo modelos sem esse ajuste (Seção 4.2, p.~13393). O artigo também investiga outros comportamentos que se agravam com RLHF, mas esse resultado não deve ser generalizado à \textit{sycophancy}. Por abordar diversas tendências comportamentais, o estudo não fornece um protocolo dedicado à capitulação factual sob contestação multi-turno.

\section{Causas ligadas ao alinhamento por preferências}

\textcite{sharma2024towards} avaliaram cinco assistentes de IA em quatro tarefas de geração de texto e verificaram que todos exibiram \textit{sycophancy}. Ao analisar os dados de preferência usados no alinhamento, os autores mostraram que julgamentos humanos e modelos de preferência às vezes favorecem respostas sicofânticas em vez das corretas, o que indica que o próprio processo de otimização por preferências contribui para o fenômeno. O foco do trabalho está na prevalência e nas causas, e não na severidade da falha entre domínios.

\section{Dinâmica em diálogos multi-turno}

\textcite{hong2025measuring} propuseram o SYCON Bench, que avalia 17 LLMs em diálogos com múltiplos turnos nos quais o usuário pressiona o modelo. O trabalho introduz duas métricas: o \textit{Turn of Flip}, que indica em qual turno o modelo muda de posição pela primeira vez, e o \textit{Number of Flip}, que conta quantas vezes ele muda de posição ao longo do diálogo. Os autores também testaram estratégias de mitigação, e uma estratégia baseada em perspectiva de terceira pessoa reduziu a \textit{sycophancy} em até 63,8\% no cenário de debate. Apesar disso, a avaliação não é centrada nas consequências das falhas entre domínios.

\section{Seletividade à correção}

\textcite{sinha2026sycobench} propôs o SycoBench-600, composto por 600 instâncias de múltipla escolha em inglês, derivadas de 272 enunciados normalizados, com três tipos de pressão e avaliação de sete assistentes. A principal contribuição é separar dois comportamentos que costumam ser medidos juntos: aceitar uma correção válida, que é desejável, e ceder a uma sugestão errada, que caracteriza a \textit{sycophancy}. O \textit{benchmark}, porém, usa questões de múltipla escolha em inglês, o que reduz a dinâmica de uma conversa livre.

\section{Síntese e lacuna}

A \autoref{tab:trabalhos} resume os trabalhos apresentados.

\begin{table}[ht]
\caption{Comparação entre os trabalhos relacionados.}
\label{tab:trabalhos}
\resizebox{\textwidth}{!}{
\begin{tabular}{|l|l|p{5.5cm}|p{5.5cm}|}
\hline
\textbf{Trabalho} & \textbf{Foco} & \textbf{Contribuição} & \textbf{Lacuna para este trabalho} \\ \hline
\textcite{perez2023discovering} & Avaliações comportamentais & Modelos maiores tendem a acompanhar posições informadas pelo usuário; a \textit{sycophancy} foi semelhante entre diferentes quantidades de etapas de RLHF. & Não avalia capitulação factual sob contestação em diálogos multi-turno. \\ \hline
\textcite{sharma2024towards} & RLHF / preferências & Cinco assistentes, quatro tarefas; preferências humanas e modelos de preferência às vezes favorecem respostas sicofânticas. & Foca prevalência e causas, não a severidade entre domínios. \\ \hline
\textcite{hong2025measuring} & Multi-turno & SYCON Bench com 17 LLMs, \textit{Turn of Flip} e \textit{Number of Flip}; mitigação em 3ª pessoa reduz até 63,8\% no debate. & Não é centrado em consequências entre domínios. \\ \hline
\textcite{sinha2026sycobench} & Seletividade & SycoBench-600: 600 instâncias sobre 272 enunciados, três tipos de pressão, sete assistentes; separa aceitar correção de ceder à sugestão errada. & Múltipla escolha em inglês; menor dinâmica de conversa livre. \\ \hline
\end{tabular}
}
\end{table}

Os quatro estudos selecionados investigam diferentes aspectos da \textit{sycophancy}, mas não organizam seus experimentos por uma classificação explícita da consequência potencial de respostas incorretas em cada cenário. Como proposta, este TCC pretende examinar capitulação e correção seletiva em cenários de diferentes níveis de criticidade. Essa delimitação se refere aos trabalhos aqui analisados; sua originalidade perante a literatura mais ampla ainda depende de revisão complementar.
